\chapter{AgenTwin Architecture}
\label{chap:architecture}

This chapter presents the research design and system architecture of AgenTwin.
The research design a Design Science Research (DSR) approach,
    mapping each DSR activity to the corresponding thesis chapter.
The system architecture then defines the layered Digital Ecosystem that
    enables agent-mediated semantic interoperability between heterogeneous Digital Twin platforms.

\section{Research Design}
\label{sec:research-design}

This research adopts the Design Science Research methodology proposed by
    Peffers et al.~\cite{peffers2007}, which provides a
    systematic framework for developing and evaluating Information Technology (IT) artifacts that
    address identified organizational problems.
DSR is particularly appropriate for this work because it
    emphasizes the creation and evaluation of
    purposeful artifacts---in this case,
    AgenTwin's Digital Ecosystem architecture and its proof-of-concept implementation.

The DSR process consists of six iterative activities:

\begin{enumerate}
    \item \textbf{Problem Identification and Motivation:}
    Addressed in Chapters~\ref{chap:introduction} and~\ref{chap:related-work},
        identifying the semantic interoperability challenge in heterogeneous DT systems.
    
    \item \textbf{Objectives of a Solution:}
    Defined in Section~\ref{sec:objectives},
        focusing on agent-mediated interoperability while preserving system autonomy.
    
    \item \textbf{Design and Development:}
    Described in Section~\ref{sec:architecture-design} and Chapter~\ref{chap:implementation},
        covering AgenTwin's Digital Ecosystem architecture and agent protocol integration.
    
    \item \textbf{Demonstration:}
    Presented in Section~\ref{sec:use-case} of Chapter~\ref{chap:evaluation},
        using a realistic cross-domain scenario involving Smart City and Energy Grid DTs.

    \item \textbf{Evaluation:}
    Detailed in Section~\ref{sec:evaluation-methodology} of Chapter~\ref{chap:evaluation}, assessing
        semantic capability,
        functional correctness and
        performance characteristics.
    
    \item \textbf{Communication:}
    Fulfilled through this thesis.
\end{enumerate}

\section{System Architecture Design}
\label{sec:architecture-design}

AgenTwin's Digital Ecosystem architecture
    depicted in Figure~\ref{fig:architecture} adopts a
    layered design inspired by
    DestinE virtual cloud layer concept~\cite{Nativi_2021}, 
    adapting its separation of concerns principles to agent-mediated Digital Twin interoperability.
Moreover, this OSI-inspired layering~\cite{day1983} provides clear separation between
    user interaction (Application),
    protocol translation (Presentation),
    conversation management (Session) and
    resource access (Domain, Integration, Platform).

\begin{figure}[htbp]
    \centering
    \includegraphics[width=\textwidth,keepaspectratio]{figures/architecture.pdf}
    \caption{
        AgenTwin architecture.
        The layered design integrates group chat patterns with MCP-based resource discovery
            to enable dynamic semantic negotiation between heterogeneous DT platforms.
    }
    \label{fig:architecture}
\end{figure}

\subsection{Application Layer}
\label{subsec:application-layer}

The application layer defines the end user that interacts with the Digital Ecosystem,
    which can be a human operator or an automated service.
A complex goal that requires assistance from multiple DT platforms is
    elaborated by the end user, then submitted to the system via agent.
Such goals may involve cross-domain queries, data aggregation, or coordinated actions
    that exceed the capabilities of any single DT platform.
Moreover, the application layer abstracts away platform-specific details,
    allowing users to focus on high-level objectives without needing to understand
    the underlying DT implementations.

\subsection{Presentation Layer}
\label{subsec:presentation-layer}

The presentation layer translates user intents into structured requests
    that can be processed by the Digital Ecosystem.
Specifically, the layer shall provide an application, service or AI agent that
    acts on behalf of the user
    to communicate with the Digital Ecosystem's agent through the A2A protocol.
This segregation allows the user application to remain agnostic of
    the underlying agent communication mechanisms,
    granting modularity and easier maintenance.
Furthermore to this point, the presentation layer may be implemented either as
    part of the Digital Ecosystem,
    integrated into the user application itself or
    provided as a third-party service,
    depending on the specific use case and system requirements.

\subsection{Session Layer}
\label{subsec:session-layer}

The session layer orchestrates multi-agent conversations that coordinate
    cross-domain interactions between heterogeneous DT platforms and user agents.
Unlike traditional orchestration approaches that rely on predetermined workflows,
    this layer enables dynamic task solving where the exact interaction sequence
    emerges from ongoing agent conversations rather than following pre-defined patterns.
This design is particularly suited for complex scenarios involving semantic interoperability,
    where the optimal coordination strategy cannot be determined a priori due to
    heterogeneous ontologies and unpredictable cross-domain dependencies.

The \textbf{Group Chat Manager} component realizes these principles through
    AutoGen's group chat pattern~\cite{wu2024autogen}, adapted for Digital Twin interoperability.
In this pattern, participating agents share the same conversation context and
    interact dynamically without strict communication order,
    enabling collaborative problem-solving that adapts to the specific characteristics
    of each cross-domain request.

The Group Chat Manager operates through three iterative steps.
First, it \textbf{dynamically selects a speaker} from available domain agents
    based on conversation context and role alignment.
This selection uses role-play prompts that consider both the current state of
    the cross-domain request and each agent's domain expertise,
    ensuring that the most appropriate agent contributes at each conversation turn.
Second, it \textbf{collects responses} from the selected agent,
    which may include semantic mappings, data queries, or requests for clarification
    from other domain agents.
Third, it \textbf{broadcasts messages} to all participating agents,
    maintaining shared context and enabling any agent to respond when their expertise is needed.

Resources are exclusively accessed through domain layer agents,
    which encapsulate platform-specific logic and knowledge.
The Group Chat Manager never directly interacts with DT platforms or their data sources;
    instead, it facilitates A2A protocol conversations among domain agents,
    which themselves use the Model Context Protocol (MCP) to
    query their respective platforms.
This protocol layering ensures clean separation between
    inter-agent coordination (A2A, session layer) and
    platform resource access (MCP, domain/integration layers).

Such design preserves platform autonomy,
    as domain agents mediate all interactions without external entities
    accessing their platforms' internals.
Digital Environment capabilities are exposed only through agent interfaces,
    enabling semantic mappings to be negotiated dynamically based on each request's context
    rather than requiring pre-coordinated schemas.
Aligned with DestinE's principles~\cite{Nativi_2021}, the session layer ensures that
    each DT platform retains full autonomy over its internal operations while
    participating in collaborative workflows mediated by agents.

Although not in this research scope, the session layer may enforce governance policies.
This includes ensuring that all interactions comply with
    security, privacy, and operational constraints
    defined by the Digital Environment administrators.
Policy enforcement occurs at the Group Chat Manager boundary, enabling centralized control
    while keeping domain agents focused on their platform-specific responsibilities.

\subsection{Domain Layer}
\label{subsec:domain-layer}

The domain layer encapsulates platform-specific agents that
    mediate interactions between the session layer and individual Digital Twin platforms.
Each \textbf{Digital Twin Agent} act as a domain specialist for a specific DT platform,
    understanding its native knowledge representation, ontology, and query capabilities.
These agents serve as semantic translators,
    exposing their platform's resources through the Model Context Protocol (MCP)
    while negotiating cross-domain mappings through the Agent-to-Agent (A2A) protocol.
To achieve semantic interoperability, besides being \textbf{domain-aware},
    each Digital Twin Agent must be capable to reason about
    semantic equivalences with other domains.
By isolating platform heterogeneity within dedicated agents,
    the domain layer enables seamless integration of diverse DT systems
    without requiring direct interoperability between platforms themselves.

\subsection{Integration Layer}
\label{subsec:integration-layer}

The integration layer provides the technical bridge between domain agents and Digital Twin platforms,
    implemented through Model Context Protocol (MCP) servers.
Each MCP server exposes a platform's capabilities through a standardized interface,
    enabling agents to discover available resources and their schemas dynamically.
This layer abstracts away platform-specific APIs and data access patterns,
    allowing agents to interact with heterogeneous DT systems through a unified protocol.

\subsection{Platform Layer}
\label{subsec:platform-layer}

The platform layer consists of individual Digital Twin systems
    that manage domain-specific digital representations of physical assets.
Each DT platform maintains its own internal data sources,
    including physical sensors, simulation models, historical datasets,
    and any other resources that contribute to its knowledge base.
DT platforms retain full autonomy over their data management,
    ensuring data integrity, security, and availability according to their own policies.
Agents access platform resources exclusively through the integration layer's MCP servers,
    preserving platform independence while enabling cross-domain coordination.
